
\section{Arity two}
\label{sec:at-ar2}

\subsection{Graphs and the local formula}

\begin{proposition}[Graph classes at arity two]
\label{prop:at-classes2}
Fix generators $x,y$ and $n\ge2$. At order $\hbar$, with $F_1=0$, the part with $n$ legs of the arity-two component along $(x,y)$ of the flow of $D$ consists of: (C2) irreducible cycles with exactly one special edge and either two vertices $S(x)$, $S(y)$ or one vertex $S^{(2)}(x,y)$ on the cycle, the other vertices from $I$; (Rd2) reducible graphs, whose subgraphs are cycles of type (C) along $x$ or $y$, of type (C2) with fewer legs, or of the arity-zero defect, joined by edges $P(\varepsilon,T)$ to trees of $I$, $S(x)$, $S(y)$, $S^{(2)}(x,y)$; (Br2) bridges; (Ct2) the counterterm trees of $D^{(2)}(x,y)$, of $D^{(1)}$ with a separate $S(y)$ or $S(x)$, and of $D^{(0)}$ with two separate vertices. The placements of each vertex are those of Lemma \ref{lem:at-place}; two vertices $S(x)$, $S(y)$ in the line placement are adjacent only through an $\AV\otimes\AV$ edge (both even: the chain $c$--$\AV$--$\AV$--$c$, a double composite line with both ends on ghost legs) or through a $\lambda\otimes\lambda$ edge (both odd, legs $\lambda$ of opposite chiralities), never when one is even and one odd (Lemma o-legs). As $\varepsilon\to0$, (Br2) with the special-edge parts of (Ct2) tends to $0$, (Rd2) to the trees of the components of $\Ob$ with lower arity or fewer legs, the $\log\varepsilon$ terms of (Ct2) cancel (Proposition \ref{prop:at-log}), (Ct2) has no pole (Proposition \ref{prop:at-poles}; $\{\int(A,A),S^{(2)}(x,y)\}=0$), and
\begin{equation}
\sum_{\mathcal C\in(C2),\,n}a_m(\mathcal C)=0\ (m\ge1),\qquad \hbar\,\Ob^{(2)}(x,y)\big|_{n\ \text{legs}}=\sum_{\mathcal C\in(C2),\,n\ \text{legs}}a_0(\mathcal C).
\label{eq:at-local2}
\end{equation}
The sign of a cycle with two vertices $S(x)$, $S(y)$ follows the arity-two component of the product of cochains (Theorem c8-L (iv)); the normalisation of Definition \ref{def:at-names} (i) is fixed by the gate G-sign (Section \ref{sec:at-engines}).
\end{proposition}

\begin{proof}
As Theorem \ref{thm:at-classes}, with Lemma \ref{lem:at-defect} for $D^{(2)}$: its order-$\hbar^0$ marked vertices are the special-edge parts of $(-1)^{|x|}\{S(x),S(y)\}_\varepsilon$ and $\{I,S^{(2)}(x,y)\}_\varepsilon$, and the marked vertices of lower arity enter with the remaining generators on separate vertices. $\{\int(A,A),S^{(2)}(x,y)\}=0$ since $\int(A,A)$ has no antifield and $S^{(2)}$ no leg $A^*$. The adjacency rules are Lemma o-legs.
\end{proof}

\begin{theorem}[The arity-two equation and its cocycle]
\label{thm:at-eq2}
Let $F_1^{(0)}$, $F_1^{(1)}$ solve the equations of arity zero and one.
\begin{itemize}
\item (i) $\mathcal R^{(2)}(x,y)$ is $\Lambda^{(0)}$-closed, of degree $-1$, Weyl weight $4-\tfrac12(\lambda_x+\lambda_y)$, position degree at most $p_x+p_y$; the arity-two equation along $(x,y)$ is solvable exactly when $[\mathcal R^{(2)}(x,y)]=0$ in $H^{-1}(\Lambda^{(0)})$.
\item (ii) Replacing $F_1^{(1)}$ by $F_1^{(1)}+\varphi$ with $\Lambda^{(0)}\varphi(x)=0$ for all $x$ replaces $\mathcal R^{(2)}(x,y)$ by
\[
\mathcal R^{(2)}(x,y)-\varphi([x,y])-(-1)^{|x|}\big(\{\varphi(x),S(y)\}+\{S(x),\varphi(y)\}\big),
\]
minus the Chevalley--Eilenberg coboundary of the one-cochain $[\varphi]$ with values in $H^{-1}(\Lambda^{(0)})$; $\Lambda^{(0)}$-exact $\varphi$ change $\mathcal R^{(2)}$ by exact terms.
\item (iii) The bracket $\{\hbar F_1^{(1)}(x),\hbar F_1^{(1)}(y)\}$ is of order $\hbar^2$ and enters at two loops; at one loop $F_1^{(1)}$ enters the arity-two equation linearly, through $F_1^{(1)}([x,y])$ and its brackets with $S(x)$, $S(y)$. $(x,y)\mapsto[\mathcal R^{(2)}(x,y)]$ is a two-cocycle of $\mathfrak{psl}(4|4)$ with values in $H^{-1}(\Lambda^{(0)})$ (arity-three component of $s(\Ob-sF_1)=0$, with $S^{(3)}=0$), and solvability through arity two for all pairs, for some choice of $F_1^{(1)}$, is the vanishing of its class in $H^2\big(\mathfrak{psl}(4|4);H^{-1}(\Lambda^{(0)})\big)$.
\item (iv) Under Hypothesis \ref{hyp:at-H}, $H^{-1}(\Lambda^{(0)})$ at the weights of the arity-two components is spanned by the classes $[S(z)]$, which form the adjoint module ($\{S(y),S(z)\}\equiv\pm S([y,z])$ modulo exact terms, Lemma c2-components); the uniform directions $\varphi(x)=t_x\,u\,S(x)$ change the class by
\begin{equation}
[\mathcal R^{(2)}(x,y)]\longmapsto[\mathcal R^{(2)}(x,y)]+u\big(t_x+t_y-t_{[x,y]}\big)[S([x,y])],
\label{eq:at-tshift}
\end{equation}
where $t_{[x,y]}S([x,y])$ means $\sum_z t_zS(z)$ over the components $z$ of $[x,y]$ in the blocks of Theorem \ref{thm:at-H}.
\end{itemize}
\end{theorem}

\begin{proof}
(i) Lemma \ref{lem:at-sformulas} (iii). (ii) Definition \ref{def:at-ob}; for exact $\varphi=\Lambda^{(0)}\psi$, $\{\Lambda^{(0)}\psi(x),S(y)\}=\Lambda^{(0)}\{\psi(x),S(y)\}$. (iii) The arity-three component of $s$ on a degree-$1$ cochain $G$ with $G^{(0)}=G^{(1)}=0$ is $\Lambda^{(0)}G^{(3)}$ minus the Chevalley--Eilenberg differential of $G^{(2)}$ for the action $y\cdot a=\pm\{S(y),a\}$, by the same forcing argument as Lemma \ref{lem:at-sformulas} (ii) applied to the arity-three identity of Lemma c2-components; in $H^{-1}$ the first term vanishes. (iv) $(-1)^{|x|}\{S(x),S(y)\}=-S([x,y])-\Lambda^{(0)}S^{(2)}(x,y)$ (Lemma c2-components), so the change in (ii) for $\varphi=tuS$ is $-u\,t_{[x,y]}S([x,y])+u(t_x+t_y)\big(S([x,y])+\Lambda^{(0)}S^{(2)}(x,y)\big)$.
\end{proof}

\begin{theorem}[What the arity-two equations fix]
\label{thm:at-fix}
Assume Hypothesis \ref{hyp:at-H}, the solvability of the arity-one equations, and write $[\mathcal R^{(2)}(x,y)]_{t=0}=u\,[S(\varrho(x,y))]$ with $\varrho(x,y)\in\mathfrak{psl}(4|4)$ of weight $\lambda_x+\lambda_y$, the class at $t_x=0$ for all $x$. The arity-two equations along all pairs are solvable exactly when
\begin{equation}
\varrho(x,y)=\big(t\cdot[x,y]\big)-[t\cdot x,y]-[x,t\cdot y]\qquad\text{for all }x,y,
\label{eq:at-rho}
\end{equation}
with $t\cdot z:=t_zz$ on the blocks, i.e. when the Chevalley--Eilenberg cocycle $\varrho$ with values in the adjoint module is the coboundary of the diagonal one-cochain $t\cdot$. Component by component:
\[
\begin{array}{l|l|l}
\text{pair} & \text{bracket} & \text{equation}\\
\hline
(\Delta,P_b) & -2P_b & \varrho=2t_\Delta P_b\\
(\Delta,\xi) & -\xi & \varrho=t_\Delta\,\xi\\
(\xi_+,\xi_-) & P_{\Gamma(\xi_+,\xi_-)|_V} & \varrho=(t_P-t_{\xi_+}-t_{\xi_-})P_{\Gamma}\\
(\xi_+,\eta_+) & \text{components }\Delta,\ M,\ R & \varrho_z=(t_z-t_{\xi_+}-t_{\eta_+})\,z,\ z\in\{\Delta,M,R\}\\
(\eta_+,\eta_-) & K_{\Gamma(\eta_+,\eta_-)|_V} & \varrho=(t_K-t_{\eta_+}-t_{\eta_-})K_{\Gamma}\\
(K_b,P_{b'}) & (b,b')\Delta+M_{b,b'} & \varrho_\Delta=(t_\Delta-t_K-t_P)(b,b')\Delta\\
 & & \varrho_M=(t_M-t_K-t_P)M_{b,b'}\\
(\xi_+,\xi'_+),\ (\xi_+,\eta_-),\ (\xi,P_b) & 0 & \varrho=0\\
(\eta_+,\eta'_+),\ (P_b,P_{b'}) & 0 & \varrho=0
\end{array}
\]
and likewise with the chiralities exchanged. Consequently:
\begin{itemize}
\item (i) $t_\Delta$, $t_M$ and $t_R$ are fixed; $t_\Delta$ is overdetermined by $(\Delta,P_b)$, $(\Delta,\xi_\pm)$ and the $\Delta$-components of $(\xi_\pm,\eta_\pm)$;
\item (ii) the solutions $(t_P,t_{\xi_\pm},t_{\eta_\pm},t_K)$ form an affine plane over the two diagonal derivations of $\mathfrak{psl}(4|4)$: the inner one, $t_z=c\lambda_z$, generated by the field redefinition $\phi^{(0)}=c\,u\,S(\Delta)$ (equivalently the finite part of the operators of Proposition c9-scheme), and the outer one, $t_z=aB_z$, the direction $G_B$ (Definition c8-B);
\item (iii) the remaining relations are consistency conditions: the agreement of the determinations of $t_\Delta$, $t_M$, $t_R$, and one linear relation among the values of $\varrho$ on $(\xi_+,\eta_+)$, $(\xi_-,\eta_-)$, $(\xi_+,\xi_-)$, $(\eta_+,\eta_-)$, $(K_b,P_{b'})$, consequences of the cocycle condition of Theorem \ref{thm:at-eq2} (iii);
\item (iv) on the commuting pairs $\varrho=0$ is automatic: $\varrho(x,y)$ depends bilinearly and equivariantly on $(x,y)$, and $\mathfrak{psl}(4|4)$ contains no element of weight $\lambda_x+\lambda_y$ of the required type ($S_+\otimes S_+$ contains no vector for $(\xi_+,\xi'_+)$ and $(\eta_+,\eta'_+)$; $S_+\otimes S_-\otimes W\otimes W$ contains no element of $\mathfrak{so}(V)\oplus\mathfrak{so}(W^\perp_{\mathbb R})\oplus\mathbb C\Delta$ for $(\xi_+,\eta_-)$; no element has weight $-3$ or $-4$ for $(\xi,P_b)$, $(P_b,P_{b'})$). The computed classes on these pairs test Hypothesis \ref{hyp:at-H} (gate G-zero).
\end{itemize}
\end{theorem}

\begin{proof}
\eqref{eq:at-tshift}: the equation along $(x,y)$ is $u[S(\varrho(x,y))]+u(t_x+t_y-t_{[x,y]})[S([x,y])]=0$, which is \eqref{eq:at-rho} componentwise. The brackets are those of Proposition c2-special, Remark c2-dilatation and Remark c7-KDelta; $[\xi_+,\eta_-]=0$ because $S_+\otimes S_-\otimes W\otimes W$ contains no element of $\mathfrak{so}(V)\oplus\mathfrak{so}(W^\perp_{\mathbb R})\oplus\mathbb C\Delta$. Solving: $(\Delta,P_b)$ gives $t_\Delta$; $(\xi_\pm,\eta_\pm)$ give $t_M-t_\Delta$, $t_R-t_\Delta$, $t_{\xi_\pm}+t_{\eta_\pm}$; $(\xi_+,\xi_-)$ gives $t_P-t_{\xi_+}-t_{\xi_-}$; $(\eta_+,\eta_-)$ gives $t_K-t_{\eta_+}-t_{\eta_-}$; $(K_b,P_{b'})$ gives $t_K+t_P$. On the six unknowns $t_P,t_{\xi_\pm},t_{\eta_\pm},t_K$ these are five equations of rank four, the sum of the equations from $(\xi_+,\eta_+)$, $(\xi_-,\eta_-)$, $(\xi_+,\xi_-)$ and $(\eta_+,\eta_-)$ being the equation from $(K_b,P_{b'})$ up to the values of $\varrho$: the homogeneous solutions are $t_z=c\lambda_z$ and $t_z=aB_z$, whose diagonal maps are derivations ($\lambda_{[x,y]}=\lambda_x+\lambda_y$; $B$ is a derivation). The inner direction: $(s\phi)^{(1)}(x)=\{\phi^{(0)},S(x)\}=c\,u\{S(\Delta),S(x)\}\equiv c\,\lambda_x\,u\,S(x)$ modulo exact terms; the outer: Proposition c8-laws (c).
\end{proof}

\subsection{The pair $(\Delta,P_b)$ and the coefficient $\theta$}

\begin{theorem}[The sector coefficients along $\Delta$ from the arity-two equation]
\label{thm:at-theta}
Let $\hbar F_1^{(1)}(\Delta)$ and $\hbar F_1^{(1)}(P_b)$ solve the arity-one equations with the quadratic parts of Definition c11-free (the particular part along $P_b$ vanishing on two legs, $\Ob^{(1)}(P_b)|_2=0$). On two legs the arity-two equation along $(\Delta,P_b)$ reads
\begin{equation}
Q\,F_1^{(2)}(\Delta,P_b)\big|_2+2u\sum_{d\in\{\perp,\lambda,V\}}\theta^\Delta_d\,S(P_b)_d=\Ob^{(2)}(\Delta,P_b)\big|_2+Q(\text{local}),
\label{eq:at-theta}
\end{equation}
independently of $\theta^{P_b}_d$. The classes $[S(P_b)_d]$ are independent in $H^{-1}(Q)$ on two legs. Hence, writing $o_d$ for the coefficient of $[S(P_b)_d]$ in $[\Ob^{(2)}(\Delta,P_b)|_2]$, the equation is solvable on two legs exactly when $[\Ob^{(2)}(\Delta,P_b)|_2]$ lies in the span of the $[S(P_b)_d]$, and then
\[
\theta^\Delta_d=\frac{o_d}{2u}\qquad(d=\perp,\lambda,V).
\]
Under Hypothesis \ref{hyp:at-H}, $o_d-2u(\theta^\Delta_d-\theta^\Delta_V)$ is independent of $d$, with the differences of Theorem \ref{thm:at-cubic}, and the components of $[\Ob^{(2)}(\Delta,P_b)|_2]$ along the second-order free symmetries of Remark \ref{rem:at-H} vanish.
\end{theorem}

\begin{proof}
In the third equation of \eqref{eq:at-eqs}, $[\Delta,P_b]=-2P_b$ and $S^{(2)}(\Delta,P_b)=0$. The sector parts of $S(\Delta)$ and of $S(P_b)$ in different sectors have no common legs, and within one sector $\{S(\Delta)_d,S(P_b)_d\}=v.S(P_b)_d=2S(P_b)_d$, the weight of $S(P_b)_d$ being $5$ (Lemma c9-SDelta (i)). The $\theta$-terms are therefore $2u\theta^\Delta_dS(P_b)_d$ from $\{F_1^{(1)}(\Delta),S(P_b)\}$, $2u\theta^{P_b}_dS(P_b)_d$ from $\{S(\Delta),F_1^{(1)}(P_b)\}$ and $-2u\theta^{P_b}_dS(P_b)_d$ from $F_1^{(1)}(-2P_b)$. The remaining quadratic parts of $F_1^{(1)}(\Delta)$ are combinations of field-counting functionals, which commute with $S(P_b)$, and $Q$-exact terms, whose brackets with $S(P_b)$ and $S(\Delta)$ are $Q$-exact. Independence: $Q$ preserves the three sectors $(\Ap,A^*|_\perp)$, $(\lambda,\lambda^*)$ and $(\AV,A^*|_V)\oplus(c,c^*)$ of the quadratic functionals of degree $-1$. In the first, the functionals of degree $-2$ and weight $5$ are $\int f(x)(A^*|_\perp,A^*|_\perp)$ with $f$ linear, whose images $\propto\int f(x)(d^*d\Ap,A^*|_\perp)$ have a nonzero position part, while $S(P_b)_\perp$ has none. In the second, the image is spanned by $Q\int(\rho(b)\lambda^*,\lambda^*)=2\int(\rho(b)\dslash\lambda,\lambda^*)=2\int(\partial_b\lambda,\lambda^*)+2\int(\rho(b\wedge\partial)\lambda,\lambda^*)$ (and position-dependent images with nonzero position part), which is not a multiple of $S(P_b)_\lambda$. In the third, the translation-invariant image is spanned by $Q\int(\iota_b\AV,c^*)$, which differs from $S(P_b)_V$ by the nonzero $\int(\iota_bd\AV,A^*|_V)$ up to sign. Hence each $[S(P_b)_d]$ is nonzero and they are independent. The last statement: $\mathcal R^{(2)}(\Delta,P_b)$ is $\Lambda^{(0)}$-closed (Theorem \ref{thm:at-eq2} (i)); by Hypothesis \ref{hyp:at-H} its class is a multiple of $[S(P_b)]$, whose two-leg part is $\sum_dS(P_b)_d$.
\end{proof}

\begin{remark}[Cross-checks for $\theta$]
\label{rem:at-thetacheck}
Theorem c9-kappa derives $\theta=-\tfrac13$ from the arity-two component of $L\circ L=0$ on the scalar bilinears, under the hypothesis that the quantisation and the local representatives exist through arity two; Corollary c11-theta does the same family by family. In the strict route this value is the prediction $o_\perp=-\tfrac23u$ for the computed $\Ob^{(2)}(\Delta,P_b)$, a cross-check (X-theta). The same $t_\Delta$ must result from $(\Delta,\xi)$ and from the $\Delta$-components of $(\xi,\eta)$ (gate G-tDelta).
\end{remark}

\subsection{The other pairs}

\begin{proposition}[The pairs along $\Sigma$ and $\Sigma^*$]
\label{prop:at-pairs}
\leavevmode
\begin{itemize}
\item (i) $(\xi,\xi')$, $\xi\in S_+\otimes W$, $\xi'\in S_-\otimes W^*$: $S^{(2)}(\xi,\xi')=\int\big(\tfrac14(\Gamma(\xi,\xi'),\Gamma(\lambda^*,\lambda^*))-\tfrac12(\xi,\lambda^*)(\xi',\lambda^*)-(\iota_{\Gamma(\xi,\xi')}A,c^*)\big)$; the two-leg part of the equation is $QF_1^{(2)}(\xi,\xi')|_2+F_1^{(1)}(P_\Gamma)|_2-\big(\{F_1^{(1)}(\xi),S(\xi')\}+\{S(\xi),F_1^{(1)}(\xi')\}\big)|_2+\{F_1^{(0)},S^{(2)}(\xi,\xi')\}|_2=\Ob^{(2)}(\xi,\xi')|_2$. The sector coefficients enter through $u\theta^{P}_dS(P_\Gamma)_d$ and through the brackets $\{S^d_2(\xi),S_2(\xi')\}$ and $\{S_2(\xi),S^d_2(\xi')\}$, whose sum over $d$ is $-S(P_\Gamma)|_2-QS^{(2)}(\xi,\xi')$ and whose individual classes in $H^{-1}(Q)$ on two legs are combinations of the $[S(P_\Gamma)_{d'}]$ given by the Fierz identities of the arrows $\lambda\to\AV\to\lambda$ and $\lambda\to\Ap\to\lambda$ (task S2-$\xi\xi$). The equation is the one-loop part of $[Q,Q]=P$.
\item (ii) $(\xi_+,\xi'_+)$: the bracket vanishes, $S^{(2)}(\xi_+,\xi'_+)$ is given by (i) with $\Gamma(\xi_+,\xi'_+)|_V=0$ and $\Gamma|_\perp\ne0$; the class of $\mathcal R^{(2)}$ must vanish, which Hypothesis \ref{hyp:at-H} makes automatic (Theorem \ref{thm:at-fix} (iv)); the rank-one case $\xi=\xi'$ enters the twisted obstruction (Remark \ref{rem:at-twist}).
\item (iii) $(\xi,\eta)$: $[\xi,\eta]$ is the vector field $x\mapsto\Gamma(\xi,\rho(x)\eta)|_V$ plus the rotation of $W^\perp_{\mathbb R}$ of \eqref{eq:c2-oddodd}, with components along $\Delta$, $\mathfrak{so}(V)$ and $\mathfrak{so}(W^\perp_{\mathbb R})$ for $(\xi_\pm,\eta_\pm)$ and zero for $(\xi_\pm,\eta_\mp)$; $S^{(2)}(\xi,\eta)$ is \eqref{eq:c2-S2} at the Killing spinor $\rho(x)\eta$. The equation is the one-loop part of $\{Q,S\}=D+M+R$; it determines $t_\Delta$, $t_M$, $t_R$ relative to $t_{\xi}+t_\eta$.
\item (iv) $(\xi,P_b)$: $[\xi,P_b]=0$, $S^{(2)}(\xi,P_b)=0$; the class of $\mathcal R^{(2)}(\xi,P_b)$ must vanish, and under Hypothesis \ref{hyp:at-H} it does, $H^{-1}(\Lambda^{(0)})$ having no class at weight $\tfrac{11}2$; the computed class tests the hypothesis. The cycles with $S(P_b)$ at the special edge or on a leg reproduce the translation of the arity-one cycles along $\xi$; the others contain the composite line along $P_b$ (gate G-xiP).
\item (v) $(\eta,\eta')$: $[\eta_+,\eta_-]=K_{\Gamma(\eta_+,\eta_-)|_V}$; evaluated at finite $T$ (Definition c8-K); the one-loop part of $\{S,S\}=K$; Table \ref{tab:at-census2} has no sector at position degree $0$, the weight $3$ requiring the position for the legs of degree $-1$.
\item (vi) $(\Delta,\xi)$: $[\Delta,\xi]=-\xi$; the one-loop part of $[D,Q]=-Q$; it fixes $t_\Delta$ again (Theorem \ref{thm:at-fix}).
\item (vii) $(K_b,P_{b'})$: at finite $T$; the one-loop part of $[K,P]=D+M$; it fixes $t_K+t_P$ and $t_M-t_\Delta$.
\end{itemize}
\end{proposition}

\begin{proof}
Theorem \ref{thm:at-fix}, Definition c2-S2 and Proposition c2-special. (iv): $S(P_b)$ commutes with every edge, vertex and counterterm; the cycles of (C2) with $S(P_b)$ at the special edge or on a leg reproduce the translation of the arity-one cycles along $\xi$, and the remaining ones carry the composite line along $P_b$.
\end{proof}

\begin{theorem}[One-loop algebra relations]
\label{thm:at-algebra}
Assume Hypothesis \ref{hyp:at-H} and the arity-one equations. For each pair $(x,y)$ the following are equivalent:
\begin{itemize}
\item (a) the arity-two equation along $(x,y)$ is solvable;
\item (b) $[\mathcal R^{(2)}(x,y)]=0$, i.e. \eqref{eq:at-rho} holds at $(x,y)$;
\item (c) with the one-loop corrected generators $S_\hbar(z):=S(z)+\hbar\,t_z\,u\,S(z)$, $z\in\{x,y,[x,y]\}$, the order-$\hbar$ part of $(-1)^{|x|}\{S_\hbar(x),S_\hbar(y)\}+S_\hbar([x,y])$ plus $u\,S(\varrho(x,y))$ is $\Lambda^{(0)}$-exact: the one-loop corrections of the generators close under the bracket of $\mathfrak{psl}(4|4)$ once the arity-two anomaly class $u[S(\varrho(x,y))]$ is included, and the equation is solvable exactly when they compensate it.
\end{itemize}
The relations are: $(\xi,\xi')$ the one-loop correction to $[Q,Q]=P$; $(\xi,\eta)$ to $\{Q,S\}=D+M+R$; $(K_b,P_{b'})$ to $[K,P]=D+M$; $(\Delta,\xi)$ to $[D,Q]=-Q$; $(\Delta,P_b)$ to $[D,P]=-2P$; $(\eta,\eta')$ to $\{S,S\}=K$. On local operators, the map $G\mapsto-\{G,-\}$ turns the action-level identity into the one-body part of the arity-two component of $L\circ L=0$ (Corollary c8-homotopy (ii)) at order $\hbar$: the commutator of the one-body parts of $\hbar L_1^{(1)}(x)$ and $\{S(y),-\}$, plus $\hbar L_1^{(1)}([x,y])$, equals a graded commutator with $\Lambda^{(0)}$ plus $-\{\hbar F_1^{(2)}(x,y),-\}$-type terms.
\end{theorem}

\begin{proof}
(a)$\Leftrightarrow$(b): Theorem \ref{thm:at-eq2} (i). (b)$\Leftrightarrow$(c): with $(-1)^{|x|}\{S(x),S(y)\}=-S([x,y])-\Lambda^{(0)}S^{(2)}(x,y)$, the order-$\hbar$ part of $(-1)^{|x|}\{S_\hbar(x),S_\hbar(y)\}+S_\hbar([x,y])$ for $S_\hbar(z)=S(z)+\hbar t_zuS(z)$ is $u(t_x+t_y-t_{[x,y]})S([x,y])$ modulo exact terms, and \eqref{eq:at-tshift}. The operator statement: Proposition c8-laws (c) and Corollary c8-homotopy (ii), the one-body terms being $-\{\hbar F_1^{(1)}(x),-\}$ (Theorem c11-classes (iii)).
\end{proof}

\begin{remark}[The twist as a cross-check]
\label{rem:at-twist}
For a rank-one $\xi=u_0\otimes z\in S_+\otimes W$, $[\xi,\xi]=0$, and the evaluation $r_\xi$ of Lemma c8-ev maps $\Ob$ to the twisted obstruction $\Ob^{(0)}+\Ob^{(1)}(\xi)-\tfrac12\Ob^{(2)}(\xi,\xi)$ and $F_1$ to $F_1^{(0)}+F_1^{(1)}(\xi)-\tfrac12F_1^{(2)}(\xi,\xi)$, which solves the one-loop twisted master equation of $S^\xi$ exactly when the equations of arities $0,1,2$ along $\xi$ hold (no components of arity $\ge3$, Lemma o-actiontrunc). The comparison with the holomorphic twist of [BK, (3.64)--(3.70)] is a cross-check only (X-twist).
\end{remark}
